\section{举不胜举}
\markboth{100. 举不胜举}{100. 举不胜举}

前面各节中已经介绍了很多种曲线，但是我们所介绍的内容，在人类关于曲线的知识宝库中不过只占了很少的一部分。因为这毕竟只是一本通俗的小册子，不可能也不应该把篇幅搞得过分庞大。为了充分利用书中最后的这几页纸，请让我们以走马看花的方式，再简单地介绍若干种其它曲线与读者见面。

（1）卡西尼卵形线。这是到两个定点的距离的乘积等于定值的点的轨迹。设两定点间的距离为 $2c$，定值为 $a^2$，曲线方程可写成
\[
  (x^2 + y^2)^2 - 2c^2(x^2 - y^2) = a^4 - c^4.
\]
特例当 $c = a$ 时得到贝努利双纽线。卡西尼卵形线是十七世纪天文学家卡西尼在试图确定地球轨道形状时发现的，卡西尼也像开普勒一样认为地球的轨道不是圆周。

（2）笛卡儿卵形线。设 $F$ 和 $F_1$ 是两个定点，$c$ 为定长，$m$ 和 $n$ 是定数，动点 $M$ 满足条件
\[
  mMF + nMF_1 = c,
\]
则 $M$ 的轨迹叫做笛卡儿卵形线，其方程可写成
\[
  (x^2 + y^2 - 2rx)^2 - l^2(x^2 + y^2) - k = 0,
\]
式中 $r$、$l$、$k$ 是一些常数，可以用 $m$、$n$、$c$ 和 $FF_1$ 表示出来。特例当 $k = 0$ 时得到帕斯卡蜗线。笛卡儿卵形线有一个重要的光学性质：可以根据媒质的折射率，选取适当的常数 $n$ 和 $m$，使得从一个定点 $F$ 发出的光线经笛卡儿卵形线折射后全部通过第二个定点 $F_1$。利用上述性质，笛卡儿选取这种曲线作为他的透镜的截面形状，用来消除透镜的球面象差。

卡西尼卵形线和笛卡儿卵形线在各有关常数取适当值时，都能成为名副其实的卵形线，即平面凸闭曲线。下面几种特殊曲线也有类似情况。

（3）孟格尔卵形线。极坐标方程是
\[
  \rho = 2a\cos^n \theta,
\]
其中 $n$ 为正整数。

（4）假开普勒卵形线。方程是
\[
  (x^2 + y^2)^2 - 2ax^3 = 0.
\]

（5）真开普勒卵形线。方程是
\[
  (x^2 + y^2 - p^2)^2 - \varepsilon^2 x^2(x^2 + y^2) = 0.
\]

下面介绍几种螺线。

（6）伽里略螺线。极坐标方程是
\[
  \rho = a\theta^2 + b\theta + c.
\]
伽里略螺线是十七世纪发现的，当时是为了考虑下面的物理问题：在地球赤道某地的上方有一个自由落体，当它随地球一起转动时，画出什么形状的曲线？答案就是伽里略螺线。

（7）抛物螺线。极坐标方程是
\[
  (\rho - l)^2 = a^2 \theta. \quad (l > 0)
\]
在数学史上，抛物螺线曾被雅各·贝努利用来作为莱布尼兹的无穷小算法应用于研究曲线性质的一个例子。

（8）连锁螺线。极坐标方程是 $\rho^2 \theta = a^2$。

（9）费马螺线。极坐标方程是 $\rho^2 = a^2 \theta$。

下面介绍几种以动物命名的曲线。

（10）蛇形线。在第 98 节箕舌线的定义中把圆 $x^2 + y^2 = ay$ 换成 $x^2 + y^2 = cx$，其余不变，就得到牛顿蛇形线。它的方程可以写成
\[
  y(x^2 + a^2) = acx.
\]

（11）蛇尾线。设 $O$ 是抛物线在其顶点处的切线上的任意一点，则此抛物线关于 $O$ 点的垂足曲线叫做蛇尾线。它的方程可以写成
\[
  x(x^2 + y^2) = y(cy - bx).
\]
特例当 $b = 0$ 时得到蔓叶线。

（12）库尔勃蚌线。方程是
\[
  x^2 y^2 + a^2(x^2 - a^2) = 0.
\]
这种曲线的形状有点像合拢的蚌壳。

（13）鸡冠线。方程是
\[
  y^2(x^2 + y^2) + 4axy^2 + a^2(3x^2 - 2y^2) - 4a^3 x + a^4 = 0.
\]

（14）蝴蝶线。方程是
\[
  y = \sqrt{2ax} + 2b\sqrt{1 - \frac{x}{2a}}.
\]

（15）蜣螂线。方程是
\[
  4(x^2 + y^2)(x^2 + y^2 + cx)^2 = a^2(x^2 - y^2).
\]

（16）肾脏线。极坐标方程是 $\rho = a\left(1 + 2\sin\frac{\theta}{2}\right)$。

再介绍几种以植物命名的曲线。

（17）梨线。方程是
\[
  x^2 y^2 - 2a^3 x + a^4 = 0.
\]

（18）四次蔓叶线。方程是
\[
  y^2 = \frac{(x^2 + 2ax + l^2)^2}{l^2 - x^2}.
\]

（19）双叶线。方程是
\[
  x^4 + y^4 - 8axy^2 = 0.
\]

（20）三叶线。方程是
\[
  x^4 + y^4 - 4ax(x^2 - y^2) = 0.
\]

（21）四叶苜蓿线。方程是
\[
  (x^2 + y^2)^3 = a^2(x^2 - y^2).
\]

最后介绍几种以杂物命名的曲线。

（22）珍珠线。这是一族曲线，方程形如
\[
  x^s(a \pm x)^r = \frac{a^{r+s}}{b^p} y^p,
\]
式中 $p$、$r$、$s$ 为自然数。其中两个重要特例是：
\begin{align*}
  x^4 + y^4 - ax^3 &= 0; \quad \text{（史留斯珍珠线）} \\
  x^4 - y^4 - 2ax^3 &= 0. \quad \text{（隆桑珍珠线）}
\end{align*}

（23）杖头线，又叫欧道克斯杖头线。方程是
\[
  x^4 - a^2(x^2 + y^2) = 0.
\]

（24）翼琴线。方程是
\[
  y^2 = \frac{(x - p)^2(a^2 - x^2)}{(x + a)^2}.
\]

（25）沙钟线。方程是
\[
  x^2 = \frac{a^2(y^2 - b^2)}{4b^2(y^2 + b^2)}.
\]

（26）荷包线。方程是
\[
  y^4 - 4(ax + b^2)y^2 + 4(a^2 + b^2)x^2 = 0, \quad b \neq 0.
\]

（27）魔鬼线。方程是
\[
  x^4 - y^4 - b^2 x^2 + a^2 y^2 = 0.
\]

（28）风车线。方程是
\[
  x^2 y^2(x^2 + y^2) = a^2(x^2 - y^2)^2.
\]

（29）8字线。方程是
\[
  x^4 - 4a^2(x^2 - y^2) = 0.
\]

（30）童衫线。方程是
\[
  y^2 = \frac{h(x + h)^2}{x}.
\]

还有若干其他已在文献中命名的曲线，举不胜举，不再赘叙。至于那些在理论研究和实际问题中遇到的无名曲线，更是千姿百态，层出不穷。无怪乎曲线这个题材引起了人们广泛的兴趣。曲线的魅力究竟在什么地方呢？是因为它们很美吗？确实，有些曲线如花似叶，婀娜多姿，它们已被艺术家们的慧眼看中，以各种美丽的图案形式，涉足于橱窗、花布、玩具、广告和书刊装帧等远离数学的地方。不过，曲线之所以迷人，关键还不在于外形美，而在于人们非常需要它，科学和生活离不开曲线。或许你不想成为诗人，但是当你在以后的工作和学习中意外地遇到各种曲线时，你也会象诗人一样赞美道：啊，曲线！美丽的曲线，有用的曲线，形形色色的曲线！
